% Generated draft: requires mathematical and source review.
\documentclass[11pt]{article}
\usepackage[margin=25mm]{geometry}
\usepackage{fontspec}
\setmainfont{Latin Modern Roman}
\usepackage{amsmath,amssymb,mathtools,unicode-math}
\setmathfont{Latin Modern Math}
\usepackage{enumitem,needspace,xurl,hyperref}
\setlength{\parindent}{0pt}
\setlength{\parskip}{5pt}
\setlength{\emergencystretch}{3em}
\newcommand{\N}{\mathbb{N}}
\newcommand{\Z}{\mathbb{Z}}
\newcommand{\Q}{\mathbb{Q}}
\newcommand{\R}{\mathbb{R}}
\newcommand{\C}{\mathbb{C}}
\newcommand{\F}{\mathbb{F}}
\newcommand{\A}{\mathbb{A}}
\newcommand{\PP}{\mathbb{P}}
\newcommand{\E}{\mathbb{E}}
\newcommand{\eps}{\varepsilon}
\newcommand{\dd}{\,\mathrm{d}}
\begin{document}
\section*{UnsolvedMath ID: 30004578}
\textbf{Source problem code:} catalogue entry
\subsection{Definitions and mathematical statement}
\textbf{Nature and scope of the question.} The supplied record poses a historiographical question concerning ancient Greek and Roman thought. It does not specify mathematical structures, axioms, or a proposition with formal quantifiers. In particular, it supplies no definitions of ``mathematical object,'' ``empirical phenomenon,'' or the sense in which such objects might be ``distinct.'' These expressions therefore remain subjects of historical interpretation rather than formally defined terms.

\textbf{Source question.} In ancient Greek and Roman thought, what does mathematics study; are mathematical objects distinct from empirical phenomena; and how do the objects of arithmetic, geometry, and applied fields such as harmonics, astronomy, statics, dynamics, mechanics, and cartography relate?

\textbf{Unresolved qualifications.} The record does not identify particular thinkers, texts, periods within ancient Greek and Roman history, or criteria for distinguishing mathematical objects from empirical phenomena. It also does not specify the kind of relationship sought among the objects of the listed disciplines. Consequently, the question cannot be converted into a unique precise mathematical proposition from the supplied evidence. Its interrogative form does not assume either a uniform ancient conception of mathematics or an affirmative answer concerning the distinctness of mathematical objects. The supplied curation describes this as a broad historiographical question, not a single falsifiable theorem, and reports that the source offers no settled answer; this assessment has not been independently verified.
\subsection{Short English statement}
The question asks how ancient Greek and Roman thinkers understood the objects of mathematics, their distinction from empirical phenomena, and their relationships across arithmetic, geometry, and the listed applied fields. The supplied record leaves the historical corpus and key concepts unspecified, so this remains a historiographical research question requiring human review.
\subsection{Sources}
\begin{itemize}
\item Local UnsolvedMath registry, supplied complete source record: permanent ID \(30004578\), title \emph{Objects and Scope of Ancient Mathematics}, source problem code \texttt{catalogue entry}, collection \emph{Oberwolfach Reports Open Problems}.
\item Supplied source citation: \emph{History of Mathematics: A Global Cultural Approach} (2021), \emph{Oberwolfach Reports}, DOI 10.4171/owr/2020/41, pp. 2021--2092. \url{https://doi.org/10.4171/owr/2020/41}. The supplied background describes Oberwolfach Reports as documenting mathematical workshops and their open-problem discussions. No individual proposer is identified; the registry gives the proposed year as 2020.
\item Supplied literature reference, labeled ``original Oberwolfach report'': \url{https://doi.org/10.4171/owr/2020/41}.
\item Literature assessment checked at \texttt{2026-08-22}: the supplied curation labels the question ``open,'' characterizes it as a broad historiographical question rather than a single falsifiable theorem, and states that the source offers no settled answer. This is supplied dated curation derived from the OWR record, not an independently verified literature assessment. Time-sensitive status and cited claims require independent verification.
\item The registry labels the statement ``exact'' and states that it is internally readable and complete and was retained byte-for-byte. These are supplied record-level assertions; no source passage has been independently checked for this draft.
\end{itemize}
\end{document}
